\chapter{Rules first, and FISTA run properly}\label{ch:fista_depth}

The review ended with a list of experiments. Before running any of them, one more thing was done first, and it is a
habit worth keeping for the rest of your research.

\section{Write the rules before you see the numbers}

Suppose you run an experiment, look at the result, and then decide what counts as success. You will nearly always find
a reading in which you succeeded. Nobody is being dishonest; with enough small choices left open, the data pick the
choices for you. The cure is to close the choices in advance.

So on 10 November a protocol was written (\code{M1_M2_PROTOCOL.md}) \emph{before} any new COMSOL score was read. It
fixes:
\begin{itemize}
\item \textbf{which data} each number is computed on, and that nothing may be chosen on COMSOL;
\item \textbf{every grid} that will be searched (heights, strengths, depths);
\item \textbf{decision rules}: what result would confirm an idea, and what result would refute it;
\item \textbf{a fingerprint}. The file's SHA-256 hash is stored with every output. If one character of the file changed
later, the hash would no longer match.
\end{itemize}
When something really had to be added later, it went into a dated \emph{amendment} at the end, written before the number
it affects. The same method was used for the 10\um\ work (Chapter~\ref{ch:ten}).

\section{The reproduction gate}

There is a second safeguard. Before any \emph{new} quantity is computed from an old model, the model must first
reproduce its \emph{old} scores. If it cannot, something in the code or the data has drifted, and every new number would
be built on sand.

\begin{code*}{m1_rescore.py}
    worst = max(abs(g['diff']) for g in gate)
\end{code*}
\begin{code*}{m1_rescore.py}
    if worst > 1e-4:
        raise SystemExit('reproduction gate failed; stopping (protocol)')
\end{code*}

\code{gate} holds, for every COMSOL pair and component, the saved score and the score from running the model again.
The result: all eight unrolled networks reproduced their saved scores with a difference of exactly 0.0, and MAGIC-UNet
to $1.5\times10^{-14}$.

\begin{keyidea}
This is why every script from here on runs from a frozen copy of the code in its own folder on the server
(\code{experiments/...}). A result that cannot be reproduced exactly from frozen code is not a result yet.
\end{keyidea}

\section{FISTA again, in one loop}

Chapter~\ref{ch:classics} built FISTA. As a reminder, it looks for sparse codes $x$ such that the dictionary, followed by
the blur, explains the measurement $b$:
\[
\min_x\ \tfrac12\,\|\,\text{mask}\,(H D x - b)\,\|^2+\lambda\,\|x\|_1 .
\]
The whole solver is this loop:

\begin{code*}{m2_fista_depth.py}
        for s in range(1, (upto or snaps[-1]) + 1):
            residual = mask * (physics.forward(y) - b)
            update = csc.shrink(y - step * physics.adjoint(residual), lam * step)
            t_next = (1 + (1 + 4 * t * t) ** .5) / 2
            y = update + (t - 1) / t_next * (update - x)
            t = t_next
            x = update
\end{code*}

Line by line:
\begin{itemize}
\item \code{residual}: how far the current guess, blurred, is from the measurement. The mask hides the padding.
\item \code{update}: one gradient step toward a better fit, then the shrink, which pushes small code values to zero.
\item \code{t_next} and \code{y}: the momentum that makes it FISTA and not ISTA. The next step starts a little beyond
the update, in the direction it was already moving.
\end{itemize}

\section{How many iterations?}

Everything before 10 November used 300 iterations. The error had been falling from 32 to 64 to 300, and nobody had
checked further. One observation makes the check cheap: the state of FISTA after $s$ iterations does not depend on how
many more you plan to do. So one run to 3{,}000 gives the answer at every depth on the way. The code simply saves the
field at chosen moments:

\begin{code*}{m2_fista_depth.py}
            elif s in snaps:
                fields[s] = synth.forward(x)
\end{code*}

The strength $\lambda$ is chosen for each depth, on 16 simulated crops:

\begin{center}
\begin{tabular}{lccc}
\toprule
$\lambda$ (relative scale) & 300 iterations & 1{,}000 & 3{,}000 \\
\midrule
almost none ($2.3\times10^{-6}$) & 40.0\% & 82.3\% & 184.2\% \\
0.001 & 36.1\% & 55.2\% & 75.5\% \\
0.01 & 24.4\% & 25.1\% & 27.4\% \\
\textbf{0.03} & \textbf{22.4\%} & \textbf{21.3\%} & \textbf{22.8\%} \\
0.1 & 24.5\% & 23.5\% & 25.3\% \\
\bottomrule
\end{tabular}
\end{center}

Look at the top row. With almost no penalty, more iterations make things \emph{worse without limit}. The solver keeps
improving its fit to the measurement, and the measurement contains noise. Fitting noise through an inverse blur
amplifies it enormously.

\begin{keyidea}[The iteration count is a second regulariser]
Stopping early is itself a way of not fitting the noise. So FISTA has two knobs that do similar jobs: $\lambda$ and the
number of iterations. ``FISTA-300'' was never ``the converged classical solution''. It was one tuned setting, and
not the best one.
\end{keyidea}

With the best $\lambda$ at each depth, on all 64 validation crops and then on COMSOL:

\begin{center}
\begin{tabular}{lccccc}
\toprule
Iterations & Validation, 64 crops & COMSOL, 15 pairs & C1 family & \code{IC_4um} & Prior band \\
\midrule
300 & 35.5\% & 30.5\% & 31.2\% & 27.6\% & 23.8\% \\
\textbf{1{,}000} & \textbf{34.0\%} & \textbf{26.7\%} & \textbf{26.6\%} & \textbf{27.2\%} & \textbf{19.3\%} \\
3{,}000 & 36.3\% & 27.4\% & 27.0\% & 29.1\% & 20.2\% \\
\bottomrule
\end{tabular}
\end{center}

The depth was fixed by the validation column \emph{before} the COMSOL column was read (amendment A2). The last column is
from Chapter~\ref{ch:limits}: the extra iterations gain in the prior band, the one that needs a prior. Up to 1{,}000 the
dictionary keeps sharpening wire edges. Beyond that it starts fitting noise.

\begin{careful}[What this did to every earlier comparison]
FISTA at its best depth is 3.8 points better than FISTA-300. Every learned network in Chapters~\ref{ch:results}
to~\ref{ch:magic_results} was compared with a baseline that was 3.8 points too weak. Against FISTA-1000:
\begin{itemize}
\item FISTA-Net (33.3\%) and CSC-LISTA (33.8\%) are not close. They are 7 points behind.
\item MAGIC-UNet run 2 (27.0\%) does not win. It ties.
\end{itemize}
\end{careful}

\section{Later: not three depths, but a curve}

You then asked a fair question: why should the best depth be exactly 300, 1{,}000 or 3{,}000? It could be 1{,}259. So
the script was extended to record the validation error every 25 iterations. It does not store the fields, only one
number each time, so it costs about 1\% more:

\begin{code*}{m2_fista_depth.py}
            if score is not None:
                if s % a.curve_every == 0:
                    field = synth.forward(x)[..., a.pad:-a.pad, a.pad:-a.pad]
                    fields[s] = float(relative_l2(field, score).sum())
\end{code*}

At 3\um\ the curve has its minimum at \textbf{925} iterations and is flat from 750 to 1{,}150
(Figure~\ref{fig:hs_fista}). On COMSOL, FISTA-925 scores 26.9\% against 26.7\% for FISTA-1000. By the rule written
before this check, \textbf{FISTA-925 is the 3\um\ baseline from now on}. It is 0.2 points weaker on COMSOL than FISTA-1000,
which had been used for the seed verdict of Chapter~\ref{ch:consistency}; Chapter~\ref{ch:consistency} shows the verdict
against both, and it is the same.

\begin{careful}[Why not the single exact step]
The curve is flat near its minimum, and the step is chosen on 64 simulated crops. Resolving it more finely than 25
steps would fit those 64 crops, not the method. And it is never chosen on COMSOL.
\end{careful}

\begin{defend}[If your professor asks: is your classical baseline fair?]
It uses the true standoff, a dictionary learned on the same simulator as the networks, and its penalty and its depth
are both chosen on simulated data from a curve recorded every 25 iterations. At 3\um\ its best depth is 925 iterations
and it scores 26.9\% on COMSOL. That is the number a learned method has to beat.
\end{defend}

\begin{defend}[If your professor asks: why did you report FISTA-300 for so long?]
Because the error was still falling at 300 and nobody ran longer. That was a mistake, found by the review and
corrected on 10 November. The early comparisons are kept in this booklet as they were, with a warning, so the
correction can be seen.
\end{defend}
